\subsection*{Variáveis e Dimensões (Tabela 2.1)}

De acordo com o modelo do misturador de manteiga de amendoim, as cinco grandezas derivadas para o experimento e suas dimensões primárias (Massa $M$, Comprimento $L$, Tempo $T$) são:

\begin{table}[h]
\centering
\begin{tabular}{lc}
\toprule

\textbf{Grandeza Derivada} & \textbf{Dimensões} \\ \midrule
Velocidade ($V$) & $L/T$ \\
Largura da lâmina ($d$) & $L$ \\
Densidade ($\rho$) & $M/L^3$ \\
Viscosidade ($\mu$) & $M/(L \cdot T)$ \\
Força de arrasto ($F_D$) & $(M \cdot L)/T^2$ \\ \bottomrule
\end{tabular}
\end{table}

Identificamos $n=5$ variáveis e $m=3$ dimensões fundamentais ($M, L, T$).

O problema envolve o movimento de uma lâmina de largura $d$ através de um fluido com propriedades $\rho$ e $\mu$ a uma velocidade $V$. O método básico de análise dimensional permite mostrar que a força $F_D$ não precisa ser estudada isoladamente contra cada variável, mas sim dentro dos blocos adimensionais $\Pi_1$ e $\Pi_2$.

%\subsubsection*{1. Variáveis e Dimensões Físicas}

%\subsubsection*{3. Grupos Adimensionais (Alvo da Prova)}
Assim, o autor demonstra, via Teorema Pi, que os dois grupos adimensionais que governam o fenômeno são:
\begin{equation*}
    \Pi_1 = \left( \frac{\mu}{\rho V d} \right), \quad \Pi_2 = \left( \frac{F_D}{\rho V^2 d^2} \right) 
\end{equation*}

\subsubsection*{1. Equação Funcional}
A dependência funcional postulada é expressa pela função:
\begin{equation*}
    F_D = F_D(V; d, \rho, \mu) \tag{2.7} 
    %F_D = f(V, d, \rho, \mu) \tag{2.7}
    \label{eq:misturador_pasta_amendoim}
\end{equation*}

\vspace{0.5cm}

\subsubsection*{2. Prova pela Aplicação do Método Básico (Eliminação Sucessiva)}

Diferente do Teorema Pi, o método básico elimina dimensões sucessivamente para limpar a equação funcional.
Partimos da Equação \ref{eq:misturador_pasta_amendoim} e seguindo os passos do método básico descritos na Seção 2.4.1:

\begin{itemize}
\item \textbf{Passo 1: Eliminar a dimensão de Massa ($M$)} \\
%Usamos a densidade $\rho$ para eliminar a massa de $F_D$ e $\mu$:
A densidade $\rho$ é a variável mais simples com massa. Dividimos $F_D$ e $\mu$ por $\rho$:
\begin{equation*}
    %\left[ \frac{F_D}{\rho} \right] = \frac{MLT^{-2}}{ML^{-3}} = L^4 T^{-2}
    \frac{F_D}{\rho} = \frac{[ML/T^2]}{[M/L^3]} = [L^4/T^2]; \quad \frac{\mu}{\rho} = \frac{[M/LT]}{[M/L^3]} = [L^2/T]
\end{equation*}
\begin{equation*}
    \left[ \frac{\mu}{\rho} \right] = \frac{ML^{-1}T^{-1}}{ML^{-3}} = L^2 T^{-1}
\end{equation*}
A nova equação funcional (eequação revisada) é: $\frac{F_D}{\rho} = f_1(V, d, \frac{\mu}{\rho})$.


\item \textbf{Passo 2: Eliminar a dimensão de Tempo ($T$)} \\
Usamos a velocidade $V$ para eliminar o tempo dos termos restantes:
%Usamos a velocidade $V [L/T]$ para cancelar o tempo das variáveis restantes:
\begin{equation*}
    %\left[ \frac{F_D}{\rho V^2} \right] = \frac{L^4 T^{-2}}{(LT^{-1})^2} = L^2
    \frac{F_D}{\rho V^2} = \frac{[L^4/T^2]}{[L^2/T^2]} = [L^2]; \quad \frac{\mu}{\rho V} = \frac{[L^2/T]}{[L/T]} = [L]
\end{equation*}
\begin{equation*}
    \left[ \frac{\mu}{\rho V} \right] = \frac{L^2 T^{-1}}{LT^{-1}} = L
\end{equation*}
A equação revisada é: $\frac{F_D}{\rho V^2} = f_2(d, \frac{\mu}{\rho V})$.

\item \textbf{Passo 3: Eliminar a dimensão de Comprimento ($L$)} \\
Usamos a largura da lâmina $d$ para tornar os termos adimensionais:
%Usamos a largura da lâmina $d [L]$ para tornar tudo adimensional:
\begin{equation*}
%    \Pi_2 = \frac{F_D}{\rho V^2 d^2} \implies [L^2] / [L^2] = 1
    \Pi_2 = \frac{F_D}{\rho V^2 d^2} \implies \frac{[L^2]}{[L^2]} = 1; \quad \Pi_1 = \frac{\mu}{\rho V d} \implies \frac{[L]}{[L]} = 1
\end{equation*}
\begin{equation*}
    \Pi_1^{-1} = \frac{\mu}{\rho V d} \implies [L] / [L] = 1
\end{equation*}
\end{itemize}

%\subsubsection*{3. Resultado Final}
%O processo de eliminação resulta na relação entre dois grupos adimensionais:
%\textbf{Conclusão:} 
Portanto, ambos os métodos chegam à mesma relação funcional adimensional:

\begin{equation*}
    \frac{F_D}{\rho V^2 d^2} = \Phi \left( \frac{\mu}{\rho V d} \right)
\end{equation*}
Os grupos encontrados são precisamente aqueles definidos da {Eq. (2.28)} do livro \cite{PMMDym2004}:
\begin{equation*}
    \Pi_1 = \left( \frac{\mu}{\rho V d} \right), \quad \Pi_2 = \left( \frac{F_D}{\rho V^2 d^2} \right)
\end{equation*}

%\subsection*{Conclusão}
A aplicação do método básico confirma que um problema complexo com cinco variáveis físicas pode ser reduzido a uma relação entre apenas dois grupos adimensionais. Isso valida a estrutura necessária para o design de experimentos do misturador, reduzindo drasticamente a quantidade de testes necessários para prever a força de arrasto.

%Motivação: No projeto de um misturador industrial, a força de arrasto ($F_D$) depende da velocidade da lâmina ($V$), de sua largura ($d$), da densidade do fluido ($\rho$) e de sua viscosidade ($\mu$). Para evitar a execução de dezenas de experimentos isolados, aplicamos a análise dimensional para encontrar os grupos que governam o fenômeno.
